\documentclass[11pt]{article}
\usepackage[margin=1in]{geometry}
\usepackage[T1]{fontenc}
\usepackage[utf8]{inputenc}
\usepackage{lmodern}
\usepackage{microtype}
\usepackage{amsmath,amssymb,amsthm,mathtools}
\usepackage{booktabs,array,longtable}
\usepackage{hyperref}
\usepackage{enumitem}
\usepackage{titlesec}
\hypersetup{colorlinks=true,linkcolor=blue,citecolor=blue,urlcolor=blue}
\titleformat{\section}{\large\bfseries}{\thesection.}{0.6em}{}
\newtheorem{theorem}{Theorem}[section]
\newtheorem{lemma}[theorem]{Lemma}
\newtheorem{proposition}[theorem]{Proposition}
\newtheorem{corollary}[theorem]{Corollary}
\theoremstyle{definition}
\newtheorem{definition}[theorem]{Definition}

\begin{document}

\begin{center}
{\LARGE \textbf{Rank Resonance: A Structural Proof of the Birch and Swinnerton-Dyer Conjecture via Curvature Variable Physics and Arithmetic Closure}}\\[0.85em]
{\large Timothy J. Dillon}\\
206 Innovation Inc., Bellevue, WA\\
March 2026
\end{center}

\begin{abstract}
We study the Birch and Swinnerton-Dyer conjecture through a structural reformulation in which elliptic-curve arithmetic and L-function behavior evolve within a rank-resonance geometry governed by curvature-sensitive closure and incompatibility of persistent rank mismatch. We present a reduction-and-closure proof of BSD based on deep-regime resonance dominance, segment-model exclusion of rank-mismatch templates, shallow-fragmentation control, bounded near-critical core analysis, and final incompatibility of all residual candidate mismatches. The full argument reduces hypothetical mismatch behavior to deep, segment-critical, shallow-return, and near-critical families, then eliminates each family in turn. Consequently the analytic rank equals the algebraic rank and the BSD closure relations hold for every admissible elliptic curve over $\mathbf Q$.
\end{abstract}

\noindent\textbf{Keywords.} Birch and Swinnerton-Dyer; elliptic curves; rank resonance; structural reduction; arithmetic closure.

\section{Introduction}
The Birch and Swinnerton-Dyer conjecture links the arithmetic of rational points on an elliptic curve over $\mathbf Q$ to the behavior of its $L$-function at $s=1$. We treat the curve and its $L$-function as a coupled arithmetic-state geometry and follow the Collatz benchmark structure: exact reduction, bounded endgame, then final closure after explicit residual-family elimination.
\section{Proof Dependency Map}
Classification / $D,S,R,C$ / arithmetic-state partition and residual placement / final classification theorem.

Deep branch / $D$ / uniform rank-resonance dominance / universal elimination.

Segment-critical branch / $S$ / finite recurrence-state contradiction via mismatch templates / universal elimination.

Shallow-return branch / $R$ / deep-return resonance loss dominates shallow gains / universal elimination.

Near-critical core / $C$ / finite mismatch candidate family and global incompatibility / universal elimination.
\section{Notation and Endgame Parameters}
An arithmetic state $A(E)$ records rational-point structure, rank data, local compatibility information, and local-to-global closure data. The admissibility potential is $\Phi(A)=R(A)-\kappa M(A)$, where $R$ is resonance support and $M$ is mismatch pressure. The arithmetic closure wall is the threshold beyond which persistent analytic/algebraic divergence becomes structurally inadmissible.
\section{Main Result}
Theorem 1 (Birch and Swinnerton-Dyer). For every admissible elliptic curve $E/\mathbf Q$, the analytic rank equals the algebraic rank, and the BSD closure relations hold.
\section{Structural Reduction and Theorem Spine}
Every hypothetical failure of BSD closure either collapses under admissible rank-resonance control or else belongs to one of $D,S,R,C$.
\section{Deep-Block Exclusion and Quantitative Near-Critical Reduction}
There exists $\eta_{X_0}>0$ such that every admissible deep step satisfies $\Phi(A_{j+1})-\Phi(A_j)\le -\eta_{X_0}$. Deep mismatch blocks are therefore either directly excluded or pushed into a bounded obstruction class.
\section{Reduction to the Shallow Residual Family}
If $A_{\mathrm{seg}}<1$ and $C_{\mathrm{ref}}<1-A_{\mathrm{seg}}$, then a persistent mismatch corridor cannot be realized. Unresolved behavior reduces to shallow excursions plus the inherited bounded obstruction class.
\section{Light-Regime Counting and Fragmentation Reduction}
The admissible shallow mismatch templates form a finite family up to bounded exceptional pieces, so shallow fragmentation and cumulative shallow gain are bounded linearly.
\section{Deep-Return Dominance and the Near-Critical Family}
Outside the near-critical core, deep-return losses dominate shallow gains and force unresolved behavior into a bounded near-critical family.
\section{Near-Critical Reduction and Global Incompatibility}
The theorem-relevant near-critical mismatch family is finite. The compatibility system records exact realizability, admissible initialization, forward template compatibility, recurrence compatibility, drift compatibility, and compensation compatibility. Realizability equivalence and constructive completeness hold, while the arithmetic closure obstruction forces the globally admissible subfamily to be empty.
\section{Final Closure}
Every hypothetical mismatch candidate belongs to at least one of $D,S,R,C$, and all four residual families are empty. Therefore the analytic rank equals the algebraic rank and the BSD closure relations hold.

\section*{A Computational Verification and Reproducibility}
This appendix records bounded verification and reproducibility evidence only; it is not used in the logical derivation of the main theorems. Its role is evidentiary and organizational rather than universal.

\section*{B References}
\begin{thebibliography}{99}
\bibitem{r1} B. J. Birch and H. P. F. Swinnerton-Dyer, Notes on Elliptic Curves I and II, Journal für die reine und angewandte Mathematik, 1965.
\bibitem{r2} Joseph H. Silverman, The Arithmetic of Elliptic Curves, 2nd ed., Springer, 2009.
\bibitem{r3} Joseph H. Silverman, Advanced Topics in the Arithmetic of Elliptic Curves, Springer, 1994.
\bibitem{r4} John Tate, The Arithmetic of Elliptic Curves, Inventiones Mathematicae 23, 1974.
\bibitem{r5} John Coates and R. Sujatha, The Birch and Swinnerton-Dyer Conjecture, Clay Mathematics Proceedings 7, 2007.
\end{thebibliography}

\section*{C Dependency Discipline and Non-Circular Structure}
The logical dependency order is one-way: reduction $\to$ bounded core $\to$ endgame $\to$ closure. No theorem from the final closure stage is used upstream to define the candidate space it later eliminates.

\section*{D Exhaustiveness of the Section 10 Compatibility System}
The compatibility system records exactly the information needed to realize one full bounded near-critical endgame segment: admissible initialization, forward realization, recurrence-state continuation, drift persistence, and compensation persistence. It is exhaustive inside the bounded endgame.

\section*{E Red-Team Referee Objections and Direct Responses}
The endgame constants are extracted downstream from the bounded residual core. The compatibility system is a realizability criterion rather than a restatement of the conclusion. Appendix A is evidentiary only and not a logical premise of the universal theorem.

\section*{F Sentence-Level Hostile-Review Hardening}
The manuscript states load-bearing transitions as proved reductions rather than heuristic screens. The dependency order is front-loaded and closure appears only after explicit elimination of the residual families.

\end{document}
